\section{Number Spiral Diagonals --- Problem 28}
\label{prob:28}
\noindent\textbf{Source:} \url{https://github.com/KO0I/projectEuler/blob/reorganize/solutions-by-language/solutions/c/number_spiral_diagonals_28.c}

\noindent\textbf{Assessment:} Closed-form idea, but no accumulated sum.

\subsection*{Mathematical goal}
Determine the total of both diagonals in a \(1001\times1001\) clockwise number spiral. Each successive odd-width ring contributes four corner values.

\subsection*{How the C code works}
Four functions, \verb|A| through \verb|D|, attempt to express the corner sequences algebraically. \verb|main| chooses a width \verb|n=5|, calls the four functions, adds the central one, and prints the result, thereby avoiding a full matrix representation.

\subsection*{Correctness and limitations}
Every helper overwrites its global \verb|sumA|, \verb|sumB|, \verb|sumC| or \verb|sumD| on each loop iteration instead of adding a sequence of terms. The width is five, not 1001. The print statement uses integer \verb|%ld| conversions with \verb|long double| arguments, which is undefined behavior. The expected total is 669171001, but this source does not compute it.


\subsection*{Correction to historical C}
{\color{red}\textbf{Correction:} All four corners of every odd-width ring are accumulated through width 1001. The old code only retained the last term and used invalid floating-point print conversions.}
